\ifdefined\gkver\else\def\gkver{student}\fi
\documentclass[\gkver]{gaokaozhenti}
\begin{document}

\chapter{2026年四川}

\begin{enumerate}
\item
%% number: 1
%% paperName: 2026年四川高考第 1 题 4 分
%% typeId: 1
%% type: 单选题
%% chapter: 运动和力
%% point: 动量和冲量
%% method: 无
%% score: 4
%% degree: 850
%% duplicateId: 0
%% body:
2026 年 2 月，我国某科创团队发布全球首款速度达到 $10 \Ums$ 的全尺寸人形机器人、该机器人体重 $75 \Ukg$；2025 年 1 月、该团队发布的四足机器人体重 $38 \Ukg$。若两款机器人均以 $10 \Ums$ 的速度同方向运动，则二者的动量大小之差为 \xzanswer{C}

\fourchoices[answer=C]
{$1850 \Ukgms$}
{$750 \Ukgms$}
{$370 \Ukgms$}
{$37 \Ukgms$}

%% answer: C
%% memo:
\memoanswer{
根据题意，由动量表达式 $p = mv$ 可得，二者的动量大小之差为

$\Delta p = m_{1}v - m_{2}v = 370 \Ukgms$

故选 C。
}
\item
%% number: 2
%% paperName: 2026年四川高考第 2 题 4 分
%% typeId: 1
%% type: 单选题
%% chapter: 直线运动
%% point: 参考系
%% method: 无
%% score: 4
%% degree: 850
%% duplicateId: 0
%% body:
小车在水平地面上做匀速直线运动。零时刻、站在小车上的甲沿与小车运动方向平行的方向抛出一个小球，乙站在小车侧方水平地面上观测到小球做直线运动直至落地。忽略空气阻力。则 \xzanswer{A}

\fourchoices[answer=A]
{乙观测到小球的运动轨迹与地面垂直}
{乙观测到小球的加速度为零}
{甲抛球方向与小车前进方向相同}
{小球相对地面的初速度不为零}

%% answer: A
%% memo:
\memoanswer{
B．根据题意可知，小球抛出后受重力作用，加速度为重力加速度，方向竖直向下，故 B 错误；

ACD．由于乙站在小车侧方水平地面上观测到小球做直线运动直至落地，可知，小球相对地面做自由落体运动，乙观测到小球的运动轨迹与地面垂直，则小球相对地面的初速度为零，即甲抛球方向与小车前进方向相反，相对小车的速度大小与小车运动的速度大小相等，故 A 正确，CD 错误。

故选 A。
}
\item
%% number: 3
%% paperName: 2026年四川高考第 3 题 4 分
%% typeId: 1
%% type: 单选题
%% chapter: 机械波
%% point: 波长频率波速
%% method: 无
%% score: 4
%% degree: 650
%% duplicateId: 0
%% body:
如图所示，空气中水平放置两端开口的圆柱形长管、管内有 $a$、$b$、$c$ 三个位置，$a$、$b$ 距离为 $1 \Ucm$，$c$ 在 $b$ 右侧。持续驱动活塞使其在 $a$、$b$ 间做周期为 $0.1 \Us$ 的简谐运动，管内形成稳定机械波。声波在空气中的传播速度取 $340 \Ums$，管足够长。则 \xzanswer{D}
\onepicture[w=7.5cm]{figs/03.png}

\fourchoices[answer=D]
{管内机械波为横波}
{管内机械波的振幅为 $1 \Ucm$}
{管内机械波的波长为 $1 \Ucm$}
{每秒有 $10$ 个完整的波经过位置 $c$}

%% answer: D
%% memo:
\memoanswer{
A．根据题意可知，管内机械波振动方向与传播方向平行，则管内机械波为纵波，故 A 错误；

B．活塞在 $ab$ 间做简谐运动，且 $ab$ 间距离为 $1 \Ucm$，则管内机械波的振幅为 $0.5 \Ucm$，故 B 错误；

C．管内机械波的波长为 $\lambda = vT = 34 \Um$，故 C 错误；

D．管内机械波的周期为 $0.1 \Us$，则每秒有 $n = \frac{t}{T} = 10$ 个完整的波经过位置 $c$，故 D 正确；

故选 D。
}
\item
%% number: 4
%% paperName: 2026年四川高考第 4 题 4 分
%% typeId: 1
%% type: 单选题
%% chapter: 曲线运动
%% point: 天体运动
%% method: 无
%% score: 4
%% degree: 650
%% duplicateId: 0
%% body:
离地球 $280$ 光年外有一恒星 TOI-561。与 TOI-561 相距约 $0.01\UAU$（日地距离为 $1\UAU$）的行星绕其公转的周期约为地球公转周期的 $\frac{1}{800}$、该行星和地球的公转均视为匀速圆周运动。则 TOI-561 与太阳的质量的比值约为 \xzanswer{B}

\fourchoices[answer=B]
{$0.16$}
{$0.64$}
{$1.6$}
{$6.4$}

%% answer: B
%% memo:
\memoanswer{
根据题意，由万有引力提供向心力有 $\frac{GMm}{r^{2}} = m\frac{4\pi^{2}}{T^{2}}r$

可得 $M = \frac{4\pi^{2}r^{3}}{GT^{2}}$

则有 $\frac{M_{\text{T}}}{M_{\text{太}}} = \frac{r_{\text{行}}^{3} \cdot T_{\text{地}}^{2}}{r_{\text{地}}^{3} \cdot T_{\text{行}}^{2}} = \left(\frac{0.01}{1}\right)^{3} \times \left(\frac{1}{\frac{1}{800}}\right)^{2} = 0.64$

故选 B。
}
\item
%% number: 5
%% paperName: 2026年四川高考第 5 题 4 分
%% typeId: 1
%% type: 单选题
%% chapter: 电场
%% point: 带电粒子的平衡
%% method: 无
%% score: 4
%% degree: 550
%% duplicateId: 0
%% body:
如图所示、金属薄板 $a$、$b$、$c$、$d$ 完全相同，用长导线和开关 $S_{1}$、$S_{2}$ 连接、$a$ 与 $b$、$c$ 与 $d$ 平行且间距相等，$S_{1}$、$S_{2}$ 均断开。$a$、$b$ 带等量异种电荷，$c$、$d$ 不带电。一质量为 $m$、带电量为 $q$ 的微粒，静止在 $a$、$b$ 之间、忽略边缘效应。则 \xzanswer{C}
\onepicture[w=6.2cm]{figs/05.png}

\fourchoices[answer=C]
{$S_{1}$ 闭合、$S_{2}$ 断开，微粒向下运动}
{$S_{1}$ 断开、$S_{2}$ 闭合，微粒向上运动}
{$S_{1}$、$S_{2}$ 同时闭合，微粒向下运动}
{$S_{1}$、$S_{2}$ 同时闭合，微粒保持静止}

%% answer: C
%% memo:
\memoanswer{
带电量为 $q$ 的微粒，静止在 $a$、$b$ 之间，说明重力与电场力平衡，电场力向上。

A．$S_{2}$ 断开，$b$ 的带电量保持不变，根据电容器极板电荷等量关系，$a$ 的带电量不变，因此 $ab$ 间电场强度不变，电场力仍等于重力，微粒保持静止，A 错误；

B．同理可得，B 错误；

CD．$S_{1}$、$S_{2}$ 同时闭合相当于 $ab$、$cd$ 两个电容器并联，电荷均匀分布在两个电容器的极板上，$ab$ 电容器带电量减小，板间场强减小，电场力小于重力，微粒向下运动，故 C 正确，D 错误。

故选 C。
}
\item
%% number: 6
%% paperName: 2026年四川高考第 6 题 4 分
%% typeId: 1
%% type: 单选题
%% chapter: 电场
%% point: 带电粒子的直线运动
%% method: 无
%% score: 4
%% degree: 450
%% duplicateId: 0
%% body:
如图所示，边长为 $l$ 的绝缘菱形支架 $EFGH$ 竖直放置，$FG$ 边固定于水平地面，$\angle EFG = 60{^\circ}$。$E$、$G$ 两点各固定一带电小球，两小球带等量异种电荷。点 $H$、$F$ 间固定一光滑绝缘直轨道，另一带电小球 $q$ 从 $H$ 点沿轨道由静止下滑至 $F$ 点。重力加速度大小为 $g$、小球均可视为点电荷。则小球 $q$ 在运动过程中 \xzanswer{A}
\onepicture[w=6.0cm]{\includesvg{figs/06}}

\fourchoices[answer=A]
{某时刻所受支持力可能为零}
{到达 $F$ 点时的速率为 $\sqrt{gl}$}
{电势能先增大后减小}
{所受电场力的冲量为零}

%% answer: A
%% memo:
\memoanswer{
A．根据菱形几何性质和等量异种电荷的电场特点可知 $HF$ 轨道电势处处相等，等量异种电荷中垂线上电场方向平行于电荷连线 $EG$，若电场力方向与重力垂直轨道的分力方向相反、大小相等，则支持力 $N = 0$，这种情况是可能存在的，故 A 正确；

B．整个过程电场力不做功，只有重力做功，由几何关系得下落高度 $h = \frac{\sqrt{3}}{2}l$

根据动能定理：$mgh = \frac{1}{2}mv^{2}$

解得 $v = \sqrt{\sqrt{3}gl}$，故 B 错误；

C．$HF$ 是等势线，整个运动过程中小球电势不变，电势能不变，故 C 错误；

D．冲量是矢量，等量异种电荷中垂线上电场方向始终不变，因此电场力方向始终不变，整个过程电场力的总冲量不为零，故 D 错误。

故选 A。
}
\item
%% number: 7
%% paperName: 2026年四川高考第 7 题 4 分
%% typeId: 1
%% type: 单选题
%% chapter: 运动和力
%% point: 传送带
%% method: 无
%% score: 4
%% degree: 350
%% duplicateId: 0
%% body:
如图所示，以恒定速率运行的传送带上有甲、乙两物块，二者与传送带相对静止，由不可伸长轻绳连接，之间无间隙。甲、乙质量均为 $m$，与传送带间的动摩擦因数分别为 $\frac{1}{8}$、$\frac{1}{4}$。某时刻、对乙施加水平向右的外力使其以恒定加速度 $\frac{g}{4}$（$g$ 为重力加速度大小）运动，经时间 $t_{1}$ 撤去外力、再经时间 $t_{2}$ 绳绷直。甲、乙均可视为质点、传送带足够长。则 \xzanswer{D}
\onepicture[w=9.0cm]{\includesvg{figs/07}}

\fourchoices[answer=D]
{$t_{1}$ 和 $t_{2}$ 满足 $t_{1} < t_{2}$}
{从撤去外力到绳绷直，因摩擦产生的热量为 $\frac{mg^{2}(t_{1} - t_{2})^{2}}{32}$}
{绳绷直后瞬间甲的动能为 $\frac{mg^{2}(t_{1} - t_{2})^{2}}{128}$}
{绳绷直以后甲、乙不会发生碰撞}

%% answer: D
%% memo:
\memoanswer{
以传送带为参考系，初始甲乙均静止。

A．施加外力过程，乙的最大速度 $v = \frac{g}{4}t_{1}$

撤去外力后，乙的加速度大小 $a_{2} = \mu_{2}g = \frac{g}{4}$

经时间 $t_{2}$ 绳绷直，此时乙还有向右的速度，说明 $a_{2}t_{2} < v$

故 $t_{2} < t_{1}$，故 A 错误；

B．绳绷直时乙的速度 $v_{2} = v - a_{2}t_{2} = \frac{g}{4}(t_{1} - t_{2})$

从撤去外力到绳绷直，乙的路程（以下所指路程均为相对传送带运动的距离）为 $x_{2} = \frac{v + v_{2}}{2}t_{2}$

因摩擦产生的热量 $Q_{2} = \frac{1}{4}mg\Delta x_{2} = \frac{1}{4}mgx_{2} = \frac{mg^{2}t_{2}(2t_{1} - t_{2})}{32}$，故 B 错误；

C．绳绷直瞬间，根据动量守恒 $mv_{2} = 2mv'$

解得 $v' = \frac{g}{8}(t_{1} - t_{2})$

此速度为相对传送带的速度，此时甲的对地速度要大于 $v'$，则其动能大于

$\frac{1}{2}mv'^{2} = \frac{mg^{2}(t_{1} - t_{2})^{2}}{128}$，故 C 错误；

D．绳绷直以后甲、乙相对传送带速度均为 $v' = \frac{g}{8}(t_{1} - t_{2})$

从绷直到乙相对传送带静止，乙的路程为

$x_{2}' = \frac{v'^{2}}{2a_{2}} = \frac{g}{32}(t_{1} - t_{2})^{2}$

从绷直到甲相对传送带静止，甲的加速度为

$a_{1} = \frac{1}{8}g$

甲的路程为

$x_{1}' = \frac{v'^{2}}{2a_{1}} = \frac{g}{16}(t_{1} - t_{2})^{2}$

从施加外力到绳绷直过程乙的路程为

$x_{2}'' = \frac{v}{2}t_{1} + \frac{v + v_{2}}{2}t_{2} = \frac{g}{8}(t_{1}^{2} - t_{2}^{2} + 2t_{1}t_{2})$

比较可知 $x_{2}' + x_{2}'' - x_{1}' > 0$

故绳绷直以后甲、乙不会发生碰撞，故 D 正确。

故选 D。
}
\item
%% number: 8
%% paperName: 2026年四川高考第 8 题 6 分
%% typeId: 2
%% type: 多选题
%% chapter: 原子和原子核
%% point: 原子核的组成
%% method: 无
%% score: 6
%% degree: 750
%% duplicateId: 0
%% body:
分离铀 238 和铀 235 常采用离心分离技术，将含有铀 238 和铀 235 两种同位素的气态六氟化铀 $\ce{^{238}UF6}$ 和 $\ce{^{235}UF6}$ 在高速转动的气体离心机中进行分离，如图所示。则 \xzanswer{BD}
\onepicture[w=4.5cm]{\includesvg{figs/08}}

\fourchoices[answer=BD]
{图中 $a$ 为 $\ce{^{235}UF6}$ 分子}
{图中 $a$ 为 $\ce{^{238}UF6}$ 分子}
{铀 238 和铀 235 的中子数相同}
{铀 238 和铀 235 的质子数相同}

%% answer: BD
%% memo:
\memoanswer{
AB．离心机高速转动时，所有分子转动的角速度 $\omega$ 相同，分子做圆周运动所需向心力 $F = m\omega^{2}r$。$\ce{^{238}UF6}$ 的分子质量大于 $\ce{^{235}UF6}$，质量更大的分子需要的向心力更大，更容易发生离心运动，向半径更大的外侧（筒壁附近）聚集，因此：外侧的 $a$ 处聚集质量更大的 $\ce{^{238}UF6}$，内侧 $b$ 处聚集质量更小的 $\ce{^{235}UF6}$，故 A 错误，B 正确。

CD．铀 238 和铀 235 是铀元素的同位素，同位素的定义是质子数相同、中子数不同的同一元素的不同核素，故 C 错误，D 正确。

故选 BD。
}
\item
%% number: 9
%% paperName: 2026年四川高考第 9 题 6 分
%% typeId: 2
%% type: 多选题
%% chapter: 光学
%% point: 光的折射
%% method: 无
%% score: 6
%% degree: 450
%% duplicateId: 0
%% body:
如图所示，在折射率为 $\sqrt{3}$、厚度为 $h$ 的长方体玻璃砖上加工一 V 形凹槽，将其制成左右对称的工件，凹槽边长 $l = \frac{\sqrt{3}}{2}h$、$\theta = 60{^\circ}$。将工件放置在水平面上，用单色平行于光沿竖直方向照射工件，观测到工件底部有多个亮度不同的区域、各区域分界线用 $a$、$b$、$c$、$d$、$e$、$f$ 表示。不考虑光的反射和干涉，所有表面均视为平面。则 \xzanswer{BC}
\onepicture[w=5.2cm]{figs/09.png}

\fourchoices[answer=BC]
{$c$、$d$ 间亮度最高}
{$a$、$b$ 间和 $e$、$f$ 间亮度最高}
{$c$、$d$ 间距为 $\frac{\sqrt{3}}{6}h$}
{$a$、$b$ 间距和 $e$、$f$ 间距均为 $\frac{\sqrt{3}}{4}h$}

%% answer: BC
%% memo:
\memoanswer{
设折射角为 $i$，根据折射定律 $n = \frac{\sin 60{^\circ}}{\sin i}$

可得 $i = 30{^\circ}$

根据几何关系可知从 V 形凹槽左侧进入的光，折射光线平行于 V 形凹槽右侧的边，同理从 V 形凹槽右侧进入的光，折射光线平行于 V 形凹槽左侧的边，如图所示

\onepicture[w=5.0cm]{figs/09_an.jpg}

AB．可知 $c$、$d$ 区域无光线进入，$a$、$b$ 和 $e$、$f$ 区域除了从上射入的光线还有折射光线进入，亮度最高，$b$、$c$ 和 $d$、$e$ 区域只有折射光线进入，故 A 错误，B 正确；

C．根据几何关系可得 $h' = h - l\sin 60{^\circ} = \frac{h}{4}$

$c$、$d$ 间距为 $l_{cd} = 2 \times h'\tan 30{^\circ} = \frac{\sqrt{3}}{6}h$，故 C 正确；

D．根据几何关系可得 $a$、$b$ 间距和 $e$、$f$ 间距均为 $l_{ab} = h\tan 30{^\circ} = \frac{\sqrt{3}}{3}h$，故 D 错误。

故选 BC。
}
\item
%% number: 10
%% paperName: 2026年四川高考第 10 题 6 分
%% typeId: 2
%% type: 多选题
%% chapter: 力物体的平衡
%% point: 多力平衡
%% method: 极值
%% score: 6
%% degree: 350
%% duplicateId: 0
%% body:
如图所示，球心为 $O$、半径为 $R$ 的半球体固定于水平地面，质量为 $m$ 的杂技演员依靠双手支撑竖直倒立在球面上。双手对球面压力的作用点的连线是与地面平行、圆心为 $O'$ 的小圆的直径，压力大小均为 $N$ 且不超过 $\frac{3}{5}mg$（$g$ 为重力加速度大小）。手与球面间动摩擦因数为 $\frac{\sqrt{3}}{3}$，最大静摩擦力与滑动摩擦力大小视为相等。设 $OO'$ 长为 $h$、则 \xzanswer{AD}
\onepicture[h=6.0cm]{\includesvg{figs/10}}

\fourchoices[answer=AD]
{$h = \frac{R}{2}$ 时，演员保持平衡状态，$N$ 大小可能为 $\frac{3}{5}mg$}
{$h = \frac{R}{2}$ 时，演员可以通过增大 $N$ 往上移动}
{$h < \frac{\sqrt{3}}{4}R$ 时，演员不可能保持平衡状态}
{演员要保持平衡状态，$N$ 的最小值为 $\frac{\sqrt{3}}{4}mg$}

%% answer: AD
%% memo:
\memoanswer{
设球心 $O$ 到单侧手的作用点连线与竖直轴 $OO'$ 的夹角为 $\theta$，由几何关系得

$\cos\theta = \frac{h}{R}$，$\sin\theta = \frac{\sqrt{R^{2} - h^{2}}}{R}$

演员受力平衡：水平分量抵消，竖直方向两个手的支持力和摩擦力的竖直分量平衡重力，整理得 $2(N\cos\theta + f\sin\theta) = mg$

最大静摩擦力满足 $f \leq f_{\max} = \mu N$

当静摩擦力取最大值时，联立可得

$N = \frac{mg}{2(\cos\theta + \frac{\sqrt{3}}{3}\sin\theta)}$

A．$h = \frac{R}{2}$ 时，代入数据解得 $N_{\min} = \frac{mg}{2}$

由于 $N \leq \frac{3}{5}mg$

可知 $N$ 的范围为 $\frac{mg}{2} \leq N \leq \frac{3}{5}mg$

$N$ 大小可能为 $\frac{3}{5}mg$，故 A 正确；

B．根据题目分析可知 $h = \frac{R}{2}$ 时，演员可以处于平衡状态，在保持平衡状态的前提下，当 $N$ 增大时，静摩擦力变小，合力始终不变，能否上升与 $N$ 无关，故 B 错误；

C．根据 $N = \frac{mg}{2(\cos\theta + \frac{\sqrt{3}}{3}\sin\theta)}$

结合题意 $N \leq \frac{3}{5}mg$

当 $N = \frac{3}{5}mg$ 时，解得平衡临界 $h \approx 0.28R < \frac{\sqrt{3}}{4}R$，故 C 错误；

D．根据 $N = \frac{mg}{2(\cos\theta + \frac{\sqrt{3}}{3}\sin\theta)}$

结合辅助角公式可得

$N = \frac{mg}{\frac{4\sqrt{3}}{3}\sin\left(\theta + 60{^\circ}\right)}$

可得 $N$ 的最小值 $N = \frac{\sqrt{3}}{4}mg$，故 D 正确。

故选 AD。
}
\item
%% number: 11
%% paperName: 2026年四川高考第 11 题 6 分
%% typeId: 4
%% type: 实验题
%% chapter: 气体性质
%% point: 验证玻意耳定律
%% method: 无
%% score: 6
%% degree: 650
%% duplicateId: 0
%% body:
某兴趣小组用注射器、挂钩式电子秤、橡胶管塞、细线等探究气体等温变化的规律，实验过程如下：
\twopicture[labelA=2026四川11a, labelB=2026四川11b, numA=1, numB=2, showlabel=false, wA=7.5cm, wB=6.0cm]
{figs/11a.png}{figs/11b.png}

\begin{enumerate}
\item
如图~\ref{2026四川11a} 所示、将细线系在活塞一端、水平固定针筒，用电子秤挂钩钩住细线、沿水平方向匀速拉动活塞。记录电子秤示数，计算出活塞与针筒内壁间摩擦力大小 $f$。
\item
用橡胶管塞塞满注射器末端小管，确认装置密封。
\item
用电子秤沿水平方向 \tkanswer{缓慢}（选填“缓慢”或“快速”）拉动活塞，活塞位于刻度 $30.0 \UmL$ 处时记录第一次电子秤示数，此后针筒内空气柱的体积 $V$ 每增加 $2 \UmL$ 记录一次示数，得到多组数据。
\item
计算针筒内空气柱压强 $p$。已知大气压强为 $p_{0}$、活塞横截面积为 $S$、由电子秤示数计算出活塞所受拉力大小 $F$，则 $p = $ \tkanswer{$p_{0} + \frac{f - F}{S}$}（用 $p_{0}$、$f$、$F$、$S$ 表示）。
\item
绘制 $p-\frac{1}{V}$ 图像如图~\ref{2026四川11b} 所示。实验表明：质量一定的空气，在温度保持不变的情况下，压强 $p$ 与体积 $V$ 成 \tkanswer{反比}（选填“正比”或“反比”）。
\end{enumerate}

%% answer:
\jdanswer{
（3）缓慢；（4）$p_{0} + \frac{f - F}{S}$；（5）反比
}

%% memo:
\memoanswer{
（3）为了保持气体的温度不变，应该用电子秤沿水平方向缓慢拉动活塞；

（4）以活塞为研究对象，根据平衡条件有 $p_{0}S + f = F + pS$

可得针筒内空气柱压强 $p = p_{0} + \frac{f - F}{S}$

（5）根据 $pV = k$ 可得 $p = k\frac{1}{V}$

由图可知 $p-\frac{1}{V}$ 图像斜率不变，即 $p$ 与 $\frac{1}{V}$ 成正比，则压强 $p$ 与体积 $V$ 成反比。
}
\item
%% number: 12
%% paperName: 2026年四川高考第 12 题 10 分
%% typeId: 4
%% type: 实验题
%% chapter: 电路
%% point: 传感器
%% method: 无
%% score: 10
%% degree: 450
%% duplicateId: 0
%% body:
现有一由柔性力敏薄膜和均压板构成的压力传感器，为探究该传感器的电阻变化规律，并用其测量压力，实验小组进行了如下实验。可用器材有：

\begin{itemize}
\item
待测压力传感器（以下简称“传感器”，空载阻值约 $600 \UO$）；
\item
直流稳压电源 $E$（输出电压 $6.00 \UV$）；
\item
定值电阻 $R_{0}$（阻值 $80.0 \UO$）；
\item
滑动变阻器 $R_{1}$（最大阻值约 $15 \UO$），滑动变阻器 $R_{2}$（最大阻值约 $1000 \UO$）；
\item
电流表（量程 $0 \sim 30 \UmA$，内阻 $10.0 \UO$）；
\item
砝码（质量为 $10 \Ug$、$20 \Ug$、$50 \Ug$、$100 \Ug$、$200 \Ug$ 各 2 个）；
\item
开关 $S_{1}$、$S_{2}$，导线若干；
\item
计算机。
\end{itemize}

\threepicture[labelA=2026四川12a, labelB=2026四川12b, labelC=2026四川12c, numA=1, numB=2, numC=3, showlabel=false, h=4.0cm]
{figs/12a.png}{figs/12b.png}{figs/12c.png}

\begin{enumerate}
\item
为探究传感器电阻随压力的变化规律，实验小组将传感器水平放置并连接实验器材，如图~\ref{2026四川12a} 所示。

①滑动变阻器应选用 \tkanswer{$R_{1}$}（选填“$R_{1}$”或“$R_{2}$”）。将滑动变阻器滑片 $P$ 置于右端，闭合开关 $S_{1}$，将开关 $S_{2}$ 拨至 $b$ 端，调节滑片 $P$ 使电流表满偏，此时 $R_{0}$ 两端的电压为 \tkanswer{$2.40$} $\UV[a]$（结果保留 3 位有效数字）。

②保持滑片 $P$ 位置不变，将开关 $S_{2}$ 拨至 $a$ 端、传感器上砝码质量从零开始每次增加 $20 \Ug$，记录多组电流表读数 $I$ 和对应的砝码质量 $m$，并绘制 $I-m$ 图像如图~\ref{2026四川12b} 所示。

③根据 $I-m$ 图像，结合图~\ref{2026四川12a} 电路分析可得，当 $m = 460 \Ug$ 时，传感器的阻值为 \tkanswer{$93.8$} $\UO[a]$（结果保留 3 位有效数字）。用相同的方法处理数据，得到传感器电阻和压力的关系。
\item
实验小组利用该传感器设计了如图~\ref{2026四川12c} 所示的压力测量电路。测量压力时，计算机采集 $c$、$d$ 间电压值，处理数据得到此时传感器的电阻值，再根据实验（1）得到的传感器电阻和压力的关系显示出压力值。若在某次测量中，$c$、$d$ 间电压为 $2.00 \UV$，此时传感器的电阻为 \tkanswer{$160$} $\UO[a]$，计算机显示的压力值为 \tkanswer{$1.62$} $\UN[a]$。（结果均保留 3 位有效数字，当地重力加速度大小为 $9.80 \Umsq$）
\end{enumerate}

%% answer:
\jdanswer{
\begin{enumerate}
	\item
	① $R_{1}$，$2.40$；③ $93.8$
	\item
	$160$，$1.62$
\end{enumerate}
}

%% memo:
\memoanswer{
（1）①根据实验原理可知，用滑动变阻器与传感器并联，认为压力传感器阻值的变化不影响并联部分电路两端的电压，所以需要滑动变阻器的阻值远小于压力传感器的阻值，所以滑动变阻器应选用最大阻值小的 $R_{1}$；

电流表满偏的电流为 $30 \UmA$，此时 $R_{0}$ 两端的电压为 $U = IR_{0} = 30 \times 10^{-3} \times 80.0 \UV = 2.40 \UV$

则并联部分电路两端的电压 $U' = IR_{0} + IR_{\mathrm{A}} = 2.70 \UV$

③当 $m = 460 \Ug$ 时，电流为 $26 \UmA$，则传感器的阻值为 $R = \frac{2.70}{26 \times 10^{-3}} \UO - 10 \UO \approx 93.8 \UO$

（2）流过传感器的电流 $I_{0} = \frac{2}{80.0} \UA = 0.025 \UA$

此时传感器的阻值为 $R = \frac{6 - 2}{0.025} \UO = 160 \UO$

当 $R = 160 \UO$ 时，电流 $I' = \frac{2.70}{160 + 10} \UA \approx 15.9 \UmA$

对比图~\ref{2026四川12b} 可知砝码的质量为 $165 \Ug$，则计算机显示的压力值为 $165 \times 10^{-3} \times 9.80 \UN \approx 1.62 \UN$
}
\item
%% number: 13
%% paperName: 2026年四川高考第 13 题 10 分
%% typeId: 6
%% type: 计算题
%% chapter: 直线运动
%% point: 匀变速直线运动
%% method: 无
%% score: 10
%% degree: 650
%% duplicateId: 0
%% body:
某款国产民用无人机已实现全自动作业。如图所示，无人机完成某次任务后开始返航，此时所在位置与降落点的水平距离为 $120 \Um$，竖直距离为 $90 \Um$。无人机先以 $10 \Ums$ 的速度沿水平直线飞行至与降落点水平距离 $50 \Um$ 处，然后沿原运动方向做匀减速直线运动至降落点正上方，随后用时 $33 \Us$ 竖直下降至降落点，返航结束。求：
\begin{enumerate}
\item
无人机沿水平方向做匀减速直线运动的加速度大小。
\item
无人机从开始返航到返航结束的位移大小和平均速度大小。
\end{enumerate}
\onepicture[w=6.5cm, align=right]{figs/13.png}

%% answer:
\jdanswer{
\begin{enumerate}
	\item
	$1 \Umsq$
	\item
	$150 \Um$，$3 \Ums$
\end{enumerate}
}

%% memo:
\memoanswer{
（1）根据匀变速直线运动公式 $v^{2} = 2ax_{0}$，代入数据有 $10^{2} = 2a \times 50$

可得匀减速直线运动的加速度大小为 $a = 1 \Umsq$

（2）无人机从开始返航到返航结束的位移大小 $x = \sqrt{h^{2} + l^{2}} = \sqrt{90^{2} + 120^{2}} = 150 \Um$

无人机匀速时需要的时间为 $t_{1} = \frac{l - x_{0}}{v} = \frac{120 - 50}{10} \Us = 7 \Us$

匀减速阶段的时间为 $t_{2} = \frac{v}{a} = \frac{10}{1} \Us = 10 \Us$

由于竖直方向运动时间为 $t_{3} = 33 \Us$

所以全程的平均速度为 $\overline{v} = \frac{x}{t_{1} + t_{2} + t_{3}} = \frac{150}{7 + 10 + 33} \Ums = 3 \Ums$
}
\item
%% number: 14
%% paperName: 2026年四川高考第 14 题 12 分
%% typeId: 6
%% type: 计算题
%% chapter: 电磁感应
%% point: 电磁感应中的匀变速
%% method: 无
%% score: 12
%% degree: 450
%% duplicateId: 0
%% body:
如图所示，两根相距 $l$ 的平行金属长导轨 $EH$、$FG$ 与金属杆 $EF$ 固定连接成 U 形框。U 形框质量为 $m$，电阻不计，静止在水平绝缘桌面上，与桌面间动摩擦因数为 $\mu$。劲度系数为 $k$ 的绝缘轻弹簧一端连接杆 $EF$ 的中点，另一端与墙壁相连，弹簧水平且处于原长。导轨上静置一质量为 $m$，电阻为 $R$ 的光滑金属杆 $JK$。空间存在方向竖直向上的匀强磁场，磁感应强度为大小为 $B$。现使杆 $JK$ 在水平向右的外力作用下做匀速直线运动。杆 $JK$ 始终与导轨垂直且接触良好，弹簧始终与杆 $EF$ 垂直且在弹性限度内，最大静摩擦力与滑动摩擦力大小视为相等，重力加速度大小为 $g$。
\begin{enumerate}
\item
求 U 形框所受最大静摩擦力的大小；
\item
若弹簧保持原长，求杆 $JK$ 做匀速直线运动速度的最大值；
\item
若杆 $JK$ 运动速度大小为 $v$，当弹簧伸长量为 $x_{0}$ 时外力功率最小，求此时 U 形框的速度大小和外力功率的最小值。
\end{enumerate}
\onepicture[w=7.5cm, align=right]{figs/14.png}

%% answer:
\jdanswer{
\begin{enumerate}
	\item
	$2\mu mg$
	\item
	$\frac{2\mu mgR}{B^{2}l^{2}}$
	\item
	$v - \frac{(kx_{0} + 2\mu mg)R}{B^{2}l^{2}}$，$(kx_{0} + 2\mu mg)v$
\end{enumerate}
}

%% memo:
\memoanswer{
（1）U 形框和导体棒的总质量为 $2m$，U 形框受到的支持力 $N = 2mg$

最大静摩擦力等于滑动摩擦力，即 $f = \mu N = 2\mu mg$

（2）当导体棒速度最大时，安培力等于最大静摩擦力，此时 U 形框即将滑动，导体棒切割磁感线产生的感应电动势 $E = Blv_{m}$

感应电流 $I = \frac{E}{R} = \frac{Blv_{m}}{R}$

安培力 $F_{\mathrm{A}} = BIl = \frac{B^{2}l^{2}v_{m}}{R}$

根据平衡条件有 $F_{\mathrm{A}} = f = 2\mu mg$

解得 $v_{m} = \frac{2\mu mgR}{B^{2}l^{2}}$

（3）设 U 形框的速度为 $u$，则导体棒与 U 形框的相对速度为 $v - u$，导体棒感应电动势 $E = Bl(v - u)$

感应电流 $I = \frac{Bl(v - u)}{R}$

导体棒受到的安培力 $F_{\mathrm{A}}' = \frac{B^{2}l^{2}(v - u)}{R}$

导体棒受到的安培力与 U 形框受到的安培力大小相等，对 U 形框受力分析有

$F_{\mathrm{A}}' = kx_{0} + 2\mu mg$

联立解得 $u = v - \frac{(kx_{0} + 2\mu mg)R}{B^{2}l^{2}}$

导体棒做匀速运动，对导体棒，根据平衡条件可得外力 $F = F_{\mathrm{A}}' = kx_{0} + 2\mu mg$

此时外力 $F$ 的功率 $P = Fv = (kx_{0} + 2\mu mg)v$
}
\item
%% number: 15
%% paperName: 2026年四川高考第 15 题 16 分
%% typeId: 6
%% type: 计算题
%% chapter: 磁场
%% point: 带电粒子在磁场中的运动
%% method: 无
%% score: 16
%% degree: 350
%% duplicateId: 0
%% body:
如图所示，纸面内建有直角坐标系 $xOy$，直线 $y = -x + d$（$d > 0$）左下为 Ⅰ 区、直线 $y = -x + 3d$ 右上为 Ⅱ 区，两条直线之间在 $y$ 轴左、右两侧分别为 Ⅲ 区和 Ⅳ 区。Ⅰ 区和 Ⅱ 区有垂直纸面的匀强磁场，磁感应强度大小均为 $B$，Ⅰ 区磁场方向垂直纸面向外，Ⅱ 区磁场方向未知；Ⅲ、Ⅳ 区无磁场。Ⅲ 区有长度为 $2d$ 且平行于 $y$ 轴的细管加速器，可使进入管内的带正电微粒具有沿 $y$ 轴正方向的恒定加速度。质量为 $m$、电荷量为 $q$（$q > 0$）的微粒甲在纸面内运动，不计微粒重力，微粒不与细管加速器发生碰撞。
\begin{enumerate}
\item
若甲在 Ⅰ 区做半径为 $2d$ 的圆周运动，求该运动的速度大小；
\item
若甲做周期性运动，每次过（$d$，$0$）时速度均沿 $x$ 轴正方向，求此周期性运动的周期；
\item
若甲在 Ⅰ 区的圆周运动轨迹半径为 $2d$ 且 Ⅱ 区磁场方向垂直纸面向外。某时刻甲与静止在（$d$，$0$）处的不带电微粒乙发生弹性正碰，碰撞过程中无电荷交换，碰撞后乙沿 $x$ 轴正方向运动。乙与甲的质量之比为 $k$。要使甲与乙再次正碰，且碰撞前甲仅进入 Ⅱ 区一次，求细管加速器位置的横坐标及 $k$ 需满足的条件。
\end{enumerate}
\onepicture[w=5.5cm, align=right]{\includesvg{figs/15}}

%% answer:
\jdanswer{
\begin{enumerate}
	\item
	$v = \frac{2qBd}{m}$
	\item
	$T = \frac{(4 + 3\pi)m}{qB}$
	\item
	$x = \frac{\pi + 2 - \sqrt{9\pi^{2} + 4}}{2\pi}d$，$k = \frac{6\pi + 1 + 2\sqrt{9\pi^{2} + 4}}{3}$
\end{enumerate}
}

%% memo:
\memoanswer{
（1）洛伦兹力提供向心力 $qvB = m\frac{v^{2}}{r}$

代入 $r = 2d$ 解得 $v = \frac{2qBd}{m}$

（2）若 Ⅱ 区域的磁场方向垂直于纸面向外，则粒子甲不满足每次经过 $(d,0)$；若 Ⅱ 区域的磁场方向垂直于纸面向内，则满足如图的回旋运动；综上，Ⅱ 区域的磁场方向垂直于纸面向内，运动轨迹如图

\onepicture[w=4.5cm]{figs/15_an1.png}

根据几何关系可知

$R = d = \frac{mv'}{qB}$

解得 $v' = \frac{qBd}{m}$

总路程

$s = 2d + 2d + 2 \times \frac{3}{4} \times 2\pi R = 4d + 3\pi d$

运动周期 $T = \frac{s}{v'} = \frac{(4 + 3\pi)m}{qB}$

（3）由题知需要在 $x$ 轴上相遇，粒子甲必定在 $(3d,0)$ 处与粒子乙相遇，则碰撞后，粒子甲反向圆周运动后竖直进入细管加速器，粒子乙正向匀速运动，画出运动图示如图

\onepicture[w=4.8cm]{figs/15_an2.png}

已知甲初速度 $v_{0} = \frac{2qBd}{m}$

甲、乙发生弹性正碰，根据动量守恒、能量守恒有 $mv_{0} = mv_{1} + kmv_{2}$，$\frac{1}{2}mv_{0}^{2} = \frac{1}{2}mv_{1}^{2} + \frac{1}{2}kmv_{2}^{2}$

联立解得 $v_{1} = -\frac{k - 1}{1 + k}v_{0}$，$v_{2} = \frac{2}{1 + k}v_{0}$

粒子甲在 Ⅰ 区域运动的半径

$R_{1} = \frac{m|v_{1}|}{qB} = \frac{2k - 2}{k + 1}d$

则细管加速器的横坐标

$x = d - R_{1} = \frac{3 - k}{k + 1}d$

粒子甲在 Ⅱ 区域运动的半径 $R_{2} = R_{1} + 2d$

根据 $R_{2} = \frac{mv_{\text{甲}2}}{qB}$

解得 $v_{\text{甲}2} = \frac{2k}{k + 1}v_{0}$

粒子甲在细管加速器中做匀加速直线运动，时间为 $t$，有 $|v_{1}| + at = v_{\text{甲}2}$

联立解得 $at = \frac{2qBd}{m} = v_{0}$

根据相遇的等时性 $\frac{1}{2}at^{2} + |v_{1}|t = 2d$，$\frac{2d}{v_{2}} = \frac{2\pi m}{qB} + t$

已知 $k > 1$，化简得 $k = \frac{6\pi + 1 + 2\sqrt{9\pi^{2} + 4}}{3}$

则细管加速器的横坐标 $x = \frac{\pi + 2 - \sqrt{9\pi^{2} + 4}}{2\pi}d$
}
\end{enumerate}

\end{document}
